% fixture: babel/ngerman-basic
% tier: T3
% engine: document
% covers: babel's ngerman (babel-german.def) in OT1: its captions (Inhaltsverzeichnis, Zusammenfassung, Abbildung, Tabelle, Beweis), \today (1.~Januar), the de-1996 patterns (dehyphn.tex: hyphenmins 2 2), \frenchspacing, \umlautlow (\lower@umlaut's x-height for \"a \"o \"u; \"e and \"i high, with \bbl@allowhyphens), the shorthands "a "o "u "s "S "z "i "e "` "' "< "> "- "= "~ "" "| "/ and ngerman's "c "f "l …, \ss
% issue: #338
\documentclass{article}
\usepackage{amsthm}
\usepackage[ngerman]{babel}
\title{Ein Titel}
\author{Autor}
\begin{document}
\maketitle
\begin{abstract}
Eine kurze Zusammenfassung.
\end{abstract}
\tableofcontents
\section{Einf"uhrung}
Die Informationstheorie besch"aftigt sich mit der Quantifizierung, Speicherung und "Ubertragung von Informationen. Sie wurde von Claude Shannon begr"undet. Gr"o"se, Stra"se, Fu"sball, "Ol und "Arger, Na"ive Ko"effizient, GRO"SE, ma"zvoll und "`AUS"' sagt er.
Ein "`Zitat"' und "<Guillemets">, Bindestrich-Wort: Mann"-Schaft, Kinder"=Garten, Ma"~Strich, Rot""Gr"un, Druck"|erzeugnis, Ein"/Aus, Schi"ffahrt und Ro"lladen.
Lange W"orter: Donaudampfschifffahrtsgesellschaftskapit"an und Rechtsschutzversicherungsgesellschaften sowie Geschwindigkeitsbegrenzungen auf Autobahnen.
\begin{figure}[h]
\caption{Eine Abbildung}
\end{figure}
\begin{table}[h]
\caption{Eine Tabelle}
\end{table}
\begin{proof}
Ein Beweis mit \"a, \"o und \"u.
\end{proof}
\end{document}
